\scanpage{12}
\deleted{\unsure{à gauche}}
sous le groupe
\[
(22)\qquad\gamma(G)=\gamma(G^b)=\{\pm1\}^S
\]
et de l'extension

$(23)$\quad
\diagram{\begin{tikzcd}[ampersand replacement=\&, row sep=small]
\&\&\&\gamma(G)\arrow[d,phantom,"\Vert"]\&\\
0\arrow[r,no head]\&\mathbb Z^S\arrow[r,"2"]\&\mathbb Z^S\arrow[r]\&\deleted{\unsure{\mu_2}}(\mathbb Z/2\mathbb Z)^S\arrow[r,no head]\&0
\end{tikzcd}}

(puiss. cartésienne, d'exposant $S$, de l'ext (14))
en considérant l'homomorphisme
\[
\mathbb Z^S\longrightarrow\mathbb Z^{\widehat{\vec A}}
\]
déduit de l'application-source
\[
\widehat{\vec A}\longrightarrow S
\]
d'où une extension
\[
\unsure{(24)|(25)}\qquad0-\mathbb Z^{\widehat{\vec A}}\longrightarrow\mathcal E(G^b)\longrightarrow\gamma(G)\longrightarrow0
\]
(ne dépendant que de $G^b$). Ceci posé,

On a
\[
(25)\qquad S_0\mathcal T_{G^b}=\mathcal E(G^b)\bigl(\underbrace{\varepsilon(G^b)}_{\varepsilon(G)}\bigr)
\]
Pour passer à $S_0\mathcal T_G$, on considère l'extension
\[
(26)\qquad0\longrightarrow\mathbb Z^{\widehat A}\;-\;\mathcal E(G)\;-\;\gamma(G)\;-\;0
\]
de $\gamma(G)=\gamma(G^b)$ par $\mathbb Z^{\widehat A}$, déduite de $\mathcal E(G^b)$ (24) par l'hom (20). On trouve
\mspage{11}
